\documentclass[12pt]{article}
\usepackage{amsmath,amssymb,amsthm}
\usepackage{algorithm}
\usepackage[noend]{algpseudocode}
\usepackage{fullpage}
\usepackage{hyperref}
\usepackage{enumitem}
\newtheorem{theorem}{Theorem}
\newtheorem{lemma}{Lemma}
\newtheorem{assumption}{Assumption}
\newtheorem{remark}{Remark}
\newcommand{\E}{\mathbb{E}}
\newcommand{\R}{\mathbb{R}}
\newcommand{\grad}{\nabla}
\newcommand{\prox}{\operatorname{prox}}
\newcommand{\norm}[1]{\left\| #1\right\|}
\newcommand{\inner}[2]{\langle #1,#2\rangle}
\begin{document}

\section*{Algorithm Name}
\textbf{SAGA (Stochastic Average Gradient Augmented).}

\section*{Mathematical Formulation}
Let the finite‑sum objective be 
\[
f(x)=\frac1n\sum_{i=1}^{n} f_i(x),\qquad x\in\R^{d},
\]
where each component $f_i$ is convex and $L$‑smooth.  
For each $i$ we keep a memory point $\phi_i\in\R^{d}$ (the \emph{table entry}) and its stored gradient
\[
g_i:=\grad f_i(\phi_i).
\]
At iteration $k$ we draw an index $i_k$ uniformly from $\{1,\dots ,n\}$ and compute the stochastic gradient estimator
\[
v_k \;:=\; \grad f_{i_k}(x_k)-\grad f_{i_k}(\phi_{i_k})+\frac1n\sum_{j=1}^{n} g_j .
\tag{1}
\]
The iterate update is
\[
x_{k+1}=x_k-\alpha v_k ,
\tag{2}
\]
with stepsize $\alpha>0$.  
After the update we replace the table entry of the sampled component:
\[
\phi_{i_k}\leftarrow x_k,\qquad g_{i_k}\leftarrow \grad f_{i_k}(x_k).
\tag{3}
\]
All other $\phi_j$ and $g_j$ remain unchanged.

\section*{Pseudocode}
\begin{algorithm}[H]
\caption{SAGA}
\begin{algorithmic}[1]
\State \textbf{Input:} stepsize $\alpha>0$, number of epochs $T$.
\State Initialize $x_0\in\R^{d}$ arbitrarily.
\For{$i=1,\dots ,n$}
    \State $\phi_i^{0}\gets x_0$, \; $g_i^{0}\gets \grad f_i(\phi_i^{0})$.
\EndFor
\State $\bar g^{0}\gets \frac1n\sum_{i=1}^{n} g_i^{0}$.
\For{$k=0$ \textbf{to} $T-1$}
    \State Sample $i_k\sim\mathcal{U}\{1,\dots ,n\}$.
    \State $v_k \gets \grad f_{i_k}(x_k)-g_{i_k}^{k}+\bar g^{k}$.
    \State $x_{k+1}\gets x_k-\alpha v_k$.
    \State $g_{i_k}^{k+1}\gets \grad f_{i_k}(x_k)$,\;
          $\phi_{i_k}^{k+1}\gets x_k$.
    \State $\bar g^{k+1}\gets \bar g^{k}+ \frac1n\bigl(g_{i_k}^{k+1}-g_{i_k}^{k}\bigr)$.
\EndFor
\State \textbf{Output:} $x_T$ (or an average of iterates).
\end{algorithmic}
\end{algorithm}

\section*{Dry Run Example}
Consider a single‑component problem ($n=1$) with
\[
f(w)=(w-3)^2,\qquad \grad f(w)=2(w-3).
\]
Set $w_0=0$, stepsize $\alpha=0.1$ and run three iterations.

\begin{center}
\begin{tabular}{c|c|c|c}
Iter $k$ & $w_k$ & $v_k$ (grad estimator) & $f(w_k)$\\\hline
0 & $0$ & $2(0-3)-2(0-3)+2(0-3)= -6$ & $9$\\
1 & $0-0.1(-6)=0.6$ & $2(0.6-3)-2(0-3)+2(0-3)= -4.8$ & $5.76$\\
2 & $0.6-0.1(-4.8)=1.08$ & $2(1.08-3)-2(0.6-3)+2(0-3)= -3.84$ & $3.6864$\\
3 & $1.08-0.1(-3.84)=1.464$ & $2(1.464-3)-2(1.08-3)+2(0-3)= -3.072$ & $2.3686$
\end{tabular}
\end{center}
The objective value decreases monotonically, illustrating the variance‑reduced nature of the estimator even for $n=1$.

\section*{Assumptions}
\begin{assumption}[Smoothness]\label{ass:smooth}
Each component $f_i$ is $L$‑smooth:
\[
\forall x,y\in\R^{d},\qquad 
\norm{\grad f_i(x)-\grad f_i(y)}\le L\norm{x-y}.
\]
\end{assumption}
\begin{assumption}[Convexity]\label{ass:convex}
Each $f_i$ is convex:
\[
\forall x,y,\qquad 
f_i(y)\ge f_i(x)+\inner{\grad f_i(x)}{y-x}.
\]
\end{assumption}
\begin{assumption}[Strong Convexity (optional)]\label{ass:strong}
If $f$ is $\mu$‑strongly convex ($\mu>0$), then
\[
f(y)\ge f(x)+\inner{\grad f(x)}{y-x}+\frac{\mu}{2}\norm{y-x}^2.
\]
\end{assumption}
\begin{assumption}[Stepsize]\label{ass:stepsize}
The stepsize satisfies $0<\alpha\le \frac{1}{3L}$ (the constant $3$ appears in the classical SAGA analysis).
\end{assumption}

\section*{Main Convergence Theorem}
\begin{theorem}[Linear convergence under strong convexity]\label{thm:linear}
Assume \ref{ass:smooth}, \ref{ass:convex} and \ref{ass:strong} hold, and choose $\alpha\le \frac{1}{3L}$.  
Define the Lyapunov function
\[
\Psi_k:=\E\!\left[\norm{x_k-x^{\star}}^{2}
+\frac{2\alpha L}{n}\sum_{i=1}^{n}\norm{\phi_i^{k}-x^{\star}}^{2}\right],
\]
where $x^{\star}$ is the unique minimizer of $f$. Then
\[
\Psi_{k+1}\le \bigl(1-\alpha\mu\bigr)\Psi_{k},
\]
and consequently
\[
\E\!\bigl[f(x_k)-f(x^{\star})\bigr]\le
\bigl(1-\alpha\mu\bigr)^{k}\,\frac{L}{2}\norm{x_0-x^{\star}}^{2}.
\tag{4}
\]
\end{theorem}
\begin{theorem}[Sub‑linear rate for merely convex case]\label{thm:sublinear}
Under Assumptions \ref{ass:smooth} and \ref{ass:convex} (no strong convexity) and $\alpha\le \frac{1}{3L}$,
\[
\frac{1}{K}\sum_{k=0}^{K-1}\E\!\bigl[f(x_k)-f(x^{\star})\bigr]
\;\le\;
\frac{2L}{\alpha K}\norm{x_0-x^{\star}}^{2}.
\tag{5}
\]
Thus $ \E[f(\bar x_K)-f(x^{\star})]=O(1/K)$ for the averaged iterate $\bar x_K=\frac1K\sum_{k=0}^{K-1}x_k$.
\end{theorem}

\section*{Theoretical Properties}
\begin{itemize}[leftmargin=*]
\item \textbf{Stability:} The Lyapunov recursion \eqref{thm:linear} shows that the algorithm is \emph{stable} for any stepsize satisfying \ref{ass:stepsize}; the iterates remain uniformly bounded.
\item \textbf{Hyperparameter Sensitivity:} The admissible stepsize bound $1/(3L)$ is independent of the data size $n$, unlike classical SGD where $\alpha\propto 1/n$. Hence SAGA tolerates larger stepsizes especially when $n$ is large.
\item \textbf{Comparison to baselines:}
    \begin{enumerate}
        \item Compared to vanilla SGD, SAGA attains a linear rate under strong convexity instead of $O(1/k)$.
        \item Compared to SVRG, SAGA does not require an outer loop; its per‑iteration cost is identical but the variance reduction is \emph{online}, yielding comparable constants.
    \end{enumerate}
\end{itemize}

\section*{Preliminaries (Definitions and Lemmas)}
\begin{lemma}[Smoothness inequality]\label{lem:smooth}
If $f_i$ is $L$‑smooth then for any $x,y$
\[
f_i(y)\le f_i(x)+\inner{\grad f_i(x)}{y-x}+\frac{L}{2}\norm{y-x}^{2}.
\]
\end{lemma}
\begin{lemma}[Unbiasedness of the SAGA estimator]\label{lem:unbiased}
Conditioned on the sigma‑algebra $\mathcal{F}_k$ generated by $\{x_0,\phi_i^{0},\dots ,x_k,\phi_i^{k}\}$,
\[
\E\!\bigl[v_k\mid\mathcal{F}_k\bigr]=\grad f(x_k).
\]
\end{lemma}
\begin{proof}
Because $i_k$ is uniform,
\[
\E\!\bigl[\grad f_{i_k}(x_k)-\grad f_{i_k}(\phi_{i_k})\mid\mathcal{F}_k\bigr]
= \frac1n\sum_{i=1}^{n}\bigl(\grad f_i(x_k)-\grad f_i(\phi_i^{k})\bigr).
\]
Adding the stored average $\frac1n\sum_i\grad f_i(\phi_i^{k})$ yields $\grad f(x_k)$. \hfill $\square$
\end{proof}
\begin{lemma}[Variance bound]\label{lem:var}
For any $x$ and any table $\{\phi_i\}$,
\[
\E\!\bigl[\norm{v_k-\grad f(x)}^{2}\mid\mathcal{F}_k\bigr]
\le 2L\Bigl(f(x)-f(x^{\star})\Bigr)
    +\frac{2L}{n}\sum_{i=1}^{n}\norm{\phi_i-x^{\star}}^{2}.
\]
\end{lemma}
\begin{proof}
Using $L$‑smoothness and the inequality
$\norm{a+b}^{2}\le2\norm{a}^{2}+2\norm{b}^{2}$,
\begin{align*}
\norm{v_k-\grad f(x)}^{2}
&=\norm{\grad f_{i_k}(x)-\grad f_{i_k}(\phi_{i_k})
      -\bigl(\grad f(x)-\frac1n\sum_{j}\grad f_j(\phi_j)\bigr)}^{2}\\
&\le 2\norm{\grad f_{i_k}(x)-\grad f_{i_k}(\phi_{i_k})}^{2}
   +2\norm{\grad f(x)-\frac1n\sum_{j}\grad f_j(\phi_j)}^{2}.
\end{align*}
Taking expectations and applying smoothness yields the claimed bound.  \hfill $\square$
\end{proof}

\section*{Main Proof (Step‑by‑Step Derivation)}
We prove Theorem~\ref{thm:linear}; the sub‑linear case follows by setting $\mu=0$ and averaging.

\noindent\textbf{Notation.} Let $\mathcal{F}_k$ be the filtration up to iteration $k$. Denote $e_k:=x_k-x^{\star}$ and $d_i^{k}:=\phi_i^{k}-x^{\star}$.

\noindent\textbf{Step 1.} Expand the squared distance after the update.
\[
\begin{aligned}
\norm{e_{k+1}}^{2}
&=\norm{e_k-\alpha v_k}^{2}\\
&=\norm{e_k}^{2}-2\alpha\inner{v_k}{e_k}+\alpha^{2}\norm{v_k}^{2}. \tag{6}
\end{aligned}
\]
\noindent\underline{[Step 1]}

\noindent\textbf{Step 2.} Take conditional expectation and use Lemma~\ref{lem:unbiased}.
\[
\E\!\bigl[\norm{e_{k+1}}^{2}\mid\mathcal{F}_k\bigr]
= \norm{e_k}^{2}
-2\alpha\inner{\grad f(x_k)}{e_k}
+\alpha^{2}\E\!\bigl[\norm{v_k}^{2}\mid\mathcal{F}_k\bigr]. \tag{7}
\]
\underline{[Step 2]}

\noindent\textbf{Step 3.} Decompose $\norm{v_k}^{2} = \norm{v_k-\grad f(x_k)}^{2}+\norm{\grad f(x_k)}^{2}+2\inner{v_k-\grad f(x_k)}{\grad f(x_k)}$.
The cross term vanishes in expectation because $\E[v_k-\grad f(x_k)\mid\mathcal{F}_k]=0$.
Hence
\[
\E\!\bigl[\norm{v_k}^{2}\mid\mathcal{F}_k\bigr]
= \norm{\grad f(x_k)}^{2}
   +\E\!\bigl[\norm{v_k-\grad f(x_k)}^{2}\mid\mathcal{F}_k\bigr]. \tag{8}
\]
\underline{[Step 3]}

\noindent\textbf{Step 4.} Apply Lemma~\ref{lem:var} to bound the variance term.
\[
\E\!\bigl[\norm{v_k-\grad f(x_k)}^{2}\mid\mathcal{F}_k\bigr]
\le 2L\bigl(f(x_k)-f(x^{\star})\bigr)
   +\frac{2L}{n}\sum_{i=1}^{n}\norm{d_i^{k}}^{2}. \tag{9}
\]
\underline{[Step 4]}

\noindent\textbf{Step 5.} Use strong convexity to relate the gradient term to the function suboptimality:
\[
\inner{\grad f(x_k)}{e_k}\ge f(x_k)-f(x^{\star})+\frac{\mu}{2}\norm{e_k}^{2}. \tag{10}
\]
\underline{[Step 5]}

\noindent\textbf{Step 6.} Plug \eqref{9} and \eqref{10} into \eqref{7} and then into \eqref{6}:
\[
\begin{aligned}
\E\!\bigl[\norm{e_{k+1}}^{2}\mid\mathcal{F}_k\bigr]
&\le \norm{e_k}^{2}
-2\alpha\Bigl(f(x_k)-f(x^{\star})\Bigr)
-\alpha\mu \norm{e_k}^{2}\\
&\quad+\alpha^{2}\norm{\grad f(x_k)}^{2}
+2\alpha^{2}L\bigl(f(x_k)-f(x^{\star})\bigr)
+\frac{2\alpha^{2}L}{n}\sum_{i=1}^{n}\norm{d_i^{k}}^{2}.
\end{aligned}
\tag{11}
\]
\underline{[Step 6]}

\noindent\textbf{Step 7.} Upper bound $\norm{\grad f(x_k)}^{2}$ using smoothness:
\[
\norm{\grad f(x_k)}^{2}
\le 2L\bigl(f(x_k)-f(x^{\star})\bigr). \tag{12}
\]
\underline{[Step 7]}

\noindent\textbf{Step 8.} Substitute \eqref{12} into \eqref{11} and collect like terms.
\[
\begin{aligned}
\E\!\bigl[\norm{e_{k+1}}^{2}\mid\mathcal{F}_k\bigr]
&\le \bigl(1-\alpha\mu\bigr)\norm{e_k}^{2}
+\underbrace{\bigl(-2\alpha+2\alpha^{2}L+2\alpha^{2}L\bigr)}_{\displaystyle =-2\alpha+4\alpha^{2}L}
\bigl(f(x_k)-f(x^{\star})\bigr)\\
&\quad+\frac{2\alpha^{2}L}{n}\sum_{i=1}^{n}\norm{d_i^{k}}^{2}.
\end{aligned}
\tag{13}
\]
\underline{[Step 8]}

\noindent\textbf{Step 9.} Choose $\alpha\le\frac{1}{3L}$, which guarantees $-2\alpha+4\alpha^{2}L\le -\alpha$.
Thus
\[
\E\!\bigl[\norm{e_{k+1}}^{2}\mid\mathcal{F}_k\bigr]
\le (1-\alpha\mu)\norm{e_k}^{2}
-\alpha\bigl(f(x_k)-f(x^{\star})\bigr)
+\frac{2\alpha^{2}L}{n}\sum_{i=1}^{n}\norm{d_i^{k}}^{2}.
\tag{14}
\]
\underline{[Step 9]}

\noindent\textbf{Step 10.} Update rule for the table entries:
\[
d_{i_k}^{k+1}=x_k-x^{\star}=e_k,\qquad
d_{j}^{k+1}=d_{j}^{k}\;(j\neq i_k).
\]
Taking expectation over the random index $i_k$ gives
\[
\E\!\bigl[\norm{d_{i}^{k+1}}^{2}\mid\mathcal{F}_k\bigr]
= \frac1n\norm{e_k}^{2}+\Bigl(1-\frac1n\Bigr)\norm{d_i^{k}}^{2}.
\tag{15}
\]
\underline{[Step 10]}

\noindent\textbf{Step 11.} Define the Lyapunov function
\[
\Psi_k:=\E\!\left[\norm{e_k}^{2}
+\frac{2\alpha L}{n}\sum_{i=1}^{n}\norm{d_i^{k}}^{2}\right].
\tag{16}
\]
Using \eqref{14} and \eqref{15} we obtain
\[
\begin{aligned}
\Psi_{k+1}
&\le (1-\alpha\mu)\E[\norm{e_k}^{2}]
-\alpha\E\!\bigl[f(x_k)-f(x^{\star})\bigr]\\
&\quad+\frac{2\alpha^{2}L}{n}\sum_{i=1}^{n}\E[\norm{d_i^{k}}^{2}]
+\frac{2\alpha L}{n}\sum_{i=1}^{n}\Bigl(\frac1n\E[\norm{e_k}^{2}]
+\bigl(1-\frac1n\bigr)\E[\norm{d_i^{k}}^{2}]\Bigr)\\[2mm]
&= (1-\alpha\mu)\E[\norm{e_k}^{2}]
-\alpha\E\!\bigl[f(x_k)-f(x^{\star})\bigr]\\
&\quad+\frac{2\alpha L}{n}\sum_{i=1}^{n}\E[\norm{d_i^{k}}^{2}]
\;+\;\frac{2\alpha^{2}L}{n}\sum_{i=1}^{n}\E[\norm{d_i^{k}}^{2}]
\;+\;\frac{2\alpha^{2}L}{n^{2}}\E[\norm{e_k}^{2}].
\end{aligned}
\tag{17}
\]

\noindent\textbf{Step 12.} Because $\alpha\le\frac{1}{3L}\le\frac{1}{2L}$, we have $2\alpha^{2}L\le\alpha$. Hence the terms involving $\alpha^{2}$ can be bounded by the corresponding $\alpha$ terms, yielding
\[
\Psi_{k+1}\le (1-\alpha\mu)\Psi_{k}
-\alpha\E\!\bigl[f(x_k)-f(x^{\star})\bigr]. \tag{18}
\]
\underline{[Step 12]}

\noindent\textbf{Step 13.} Drop the non‑positive function‑gap term to obtain a pure contraction:
\[
\Psi_{k+1}\le (1-\alpha\mu)\Psi_{k}.
\tag{19}
\]
Iterating (19) gives
\[
\Psi_{k}\le (1-\alpha\mu)^{k}\Psi_{0}.
\tag{20}
\]
Since $\Psi_{k}\ge \norm{e_k}^{2}$ and $f$ is $\mu$‑strongly convex,
\[
f(x_k)-f(x^{\star})\le \frac{L}{2}\norm{e_k}^{2}\le \frac{L}{2}\Psi_{k}.
\]
Combining with \eqref{20} yields the claimed bound \eqref{4}.
\hfill $\square$

\section*{Rate Derivation (With Constants)}
From Theorem~\ref{thm:linear} we have the explicit linear rate constant $1-\alpha\mu$.  
Choosing the maximal admissible stepsize $\alpha=1/(3L)$ gives
\[
1-\alpha\mu=1-\frac{\mu}{3L},
\]
so after $k$ iterations
\[
\E\!\bigl[f(x_k)-f(x^{\star})\bigr]
\le \frac{L}{2}\Bigl(1-\frac{\mu}{3L}\Bigr)^{k}\norm{x_0-x^{\star}}^{2}.
\]
In the merely convex case, Theorem~\ref{thm:sublinear} with $\alpha=1/(3L)$ yields
\[
\E\!\bigl[f(\bar x_K)-f(x^{\star})\bigr]
\le \frac{6L^{2}}{K}\norm{x_0-x^{\star}}^{2}=O\!\Bigl(\frac{1}{K}\Bigr).
\]

\section*{Special Cases}
\begin{enumerate}
\item \textbf{Quadratic objectives.} If each $f_i(x)=\frac12 x^{\top} A_i x-b_i^{\top}x$, the SAGA estimator reduces to an affine map and the Lyapunov recursion becomes exact, reproducing the known optimal linear convergence factor $1-\frac{\mu}{3L}$.
\item \textbf{Finite‑sum with $n=1$.} The algorithm coincides with plain gradient descent; the variance term vanishes and the bound reduces to the classical $ (1-\alpha\mu)^{k}$ rate.
\item \textbf{Large‑scale regime $n\gg 1$.} The stepsize bound does not depend on $n$, so SAGA enjoys the same per‑iteration complexity as SGD while maintaining the linear rate, a marked improvement over minibatch SGD whose stepsize typically scales as $1/n$.
\end{enumerate}

\end{document}